\section{Angelic Will, Love, Creation, and Elevation to Grace}\label{sec:st-will-grace}

When a blow descends toward a man's head, his hand flashes upward before he has time to deliberate, exposing bone and flesh to injury to save the whole body. In \ST{I, q.~60, a.~5}, St.~Thomas Aquinas reads in that uncalculated reflex the deepest law of created being: because every part belongs naturally more to its whole than to itself, an angel---and a human being---naturally loves God more than itself. Questions~59 through~62 of the \emph{Prima Pars} trace the full arc of angelic appetite and destiny, from pure spiritual free-will unclouded by bodily passion, through creation out of nothing at the dawn of time, to the single act of charity by which the faithful angels merited and entered eternal vision.

\subsection{The Will and Free Choice of the Angels (\ST{I, q.~59})}

Every creature inclines toward its proper good according to its mode of knowing (\ST{I, q.~59, a.~1}). Things without knowledge tend toward their end by natural appetite alone; animals endowed with sense pursue particular sensible goods through sensitive appetite; intellectual beings, grasping the universal notion of goodness itself, tend toward good through intellective appetite---the \emph{will}. Because the angelic mind bears the image of the Trinity in memory, understanding, and will (St.~Augustine, \emph{De Trinitate}~X.11--12), will necessarily belongs to the angels. Angels do not possess an appetite higher than will, since intuitive intellect and discursive reason share the same formal object of universal goodness (\ST{I, q.~59, a.~1, ad~1}); nor does the name ``appetite'' confine the will to chasing absent goods, for it equally embraces joyful rest in goods already possessed (\ST{I, q.~59, a.~1, ad~2--3}).

In every angel, the will is a distinct faculty, really different from both its essence and its intellect (\ST{I, q.~59, a.~2}). A created essence is bounded within itself, so any inclination toward an extrinsic good requires a superadded power; only in God, whose infinite essence contains universal goodness, is will identical with essence. Will and intellect likewise differ in their opposite directions of intentionality: intellect draws external things inward so that the known exists inside the knower under the aspect of the true, whereas will goes outward toward things as they exist in themselves under the aspect of the good (\ST{I, q.~59, a.~2, ad~2--3}). As in Question~54, Franciscan theology diverges from Aquinas's real distinction of faculties, distinguishing the angelic will from essence and intellect formally or virtually.

From intellect and will flows free-will (\emph{liberum arbitrium}) in a higher perfection than in human beings (\ST{I, q.~59, a.~3}). An arrow flies toward its target solely as driven by an archer, without judgment of its own. A sheep flees a wolf by a sensory judgment implanted by nature, not chosen freely. Only an intellectual agent, comprehending the common nature of goodness and weighing particular goods in its light, acts by free judgment---confirming the early patristic teaching of St.~Irenaeus (\emph{Adversus haereses}~IV.37.1) and St.~John of Damascus (\emph{De fide orthodoxa}~II.3) that God placed the power of choice in angels as rational beings. Adapting Aristotle's account of human choice (\emph{Nicomachean Ethics}~III.3), Aquinas notes that angelic choice requires no anxious inquiry or deliberation of counsel, proceeding instead ``by the sudden acceptance of truth'' in practical matters (\ST{I, q.~59, a.~3, ad~1}); while absence of coercion admits no degrees, free-will itself stands more nobly in higher angels according to the clarity of their intellectual judgment (\ST{I, q.~59, a.~3, ad~2--3}).

Because angels lack bodies and sensory organs, they possess neither an irascible nor a concupiscible appetite, only an undivided will (\ST{I, q.~59, a.~4}). Sensitive appetite splits into concupiscible and irascible powers because sensible goods confront animals now as simply pleasurable, now as arduous to win or avoid; the will tends to good under its universal formality and therefore remains one undivided faculty. When Scripture speaks of love, joy, anger, or fear in angels, or when Pseudo-Dionysius speaks of ``unreasonable fury and wild concupiscence'' in demons (\emph{De divinis nominibus}~IV), those words signify simple acts of will or likeness of effect, not bodily passions (St.~Augustine, \emph{City of God}~IX; \ST{I, q.~59, a.~4, ad~1--2}). Charity and hope reside directly in the will even in men, while temperance and fortitude belong to angels purely as moderate conformity to God's will and steadfast execution of His commands (\ST{I, q.~59, a.~4, ad~3}).

\angeldossier{\ST{I, q.~59, aa.~1--4} (Will of the angels: existence of will, distinction from essence and intellect, free-will, and absence of irascible and concupiscible appetites).}{Church teaching (aa.~1, 3); Common teaching (a.~4); Thomist position (a.~2).}{St.~Augustine, \emph{De Trinitate}~X.11--12 and \emph{De civitate Dei}~IX; Pseudo-Dionysius, \emph{De divinis nominibus}~IV; Aristotle, \emph{De anima}~II--III and \emph{Nicomachean Ethics}~III.3; St.~Irenaeus, \emph{Adversus haereses}~IV.37.1; St.~John of Damascus, \emph{De fide orthodoxa}~II.3; Peter Lombard, \emph{Sententiae}~II, d.~3.}{Franciscan theologians distinguish angelic will from essence and intellect formally or virtually rather than by a Thomist real distinction of faculties (a.~2); none recorded on aa.~1, 3--4.}

\subsection{Natural and Elective Love in the Angels (\ST{I, q.~60})}

Because love follows knowledge---``nothing is loved except it be first known'' (St.~Augustine, \emph{De Trinitate}~X.1--2)---the twofold character of angelic life produces two loves in the angelic will: \emph{natural love} (\emph{dilectio naturalis}) and \emph{love of choice} (\emph{dilectio electiva}). First, natural love necessarily dwells in every angel (\ST{I, q.~60, a.~1}). What is prior is preserved in what is posterior; since nature precedes intellect and every created nature bears an inborn inclination to its proper good, an intellectual nature retains a natural inclination inside its will. Implanted by the Author of nature, this natural love is always upright in its own order, yet remains distinct from---and elevated by---supernatural charity (\ST{I, q.~60, a.~1, ad~1--3}).

Second, angels also exercise love of choice (\ST{I, q.~60, a.~2}), by which alone a rational creature merits or demerits. Here Aquinas marks a striking asymmetry between intellect and will. In the intellect, moving discursively from principles to conclusions betrays the weakness of human light, so angels know without discursion; in the will, however, loving the end naturally while freely choosing means ordered to that end involves no imperfection at all, for the end stands in the will just as first principles stand in the intellect (Aristotle, \emph{Physics}~II).

Third, an angel loves himself with both natural love and love of choice (\ST{I, q.~60, a.~3}). Following Aristotle's observation that ``love for others comes of love for oneself'' (\emph{Nicomachean Ethics}~IX.8), Aquinas distinguishes \emph{love of friendship} (\emph{amor amicitiae}, wishing good to a subsisting person for that person's own sake) from \emph{love of concupiscence} (\emph{amor concupiscentiae}, wishing an accidental or inherent good for someone). An angel loves himself naturally by willing his own subsisting perfection, and electively by freely choosing particular goods beneficial to himself; as an immanent act, love reflects back upon the lover just as knowledge does, and unity stands even deeper than union (\ST{I, q.~60, a.~3, ad~2--3}).

Fourth, one angel naturally loves another angel as himself (\ST{I, q.~60, a.~4}). Even irrational creatures show natural affinity for their kind: ``Every beast loveth its like'' (Ecclus~13:19). Insofar as two angels share a common angelic nature, each naturally wishes well to the other as to himself---``as himself'' signifying likeness of benevolence rather than equal intensity, since numerical unity with oneself surpasses generic unity with another (\ST{I, q.~60, a.~4, ad~2}). Even fallen demons cannot erase this natural affection for good angels insofar as both share angelic nature, though demons hate the holy angels insofar as righteousness clashes with their perverse wills (\ST{I, q.~60, a.~4, ad~3}).

Question~60 culminates in a thesis of extraordinary theological power: by natural love alone, an angel loves God before himself and more intensely than himself (\ST{I, q.~60, a.~5}). Some medieval theologians had argued the opposite---conceding that an angel naturally wishes God a greater good (namely, to be God), yet claiming that absolutely speaking a creature loves itself more intensely than God until supernatural charity reverses that natural selfishness. Aquinas dismantles that view from the natural order of part and whole:
\begin{quotation}
For we observe that the part naturally exposes itself in order to safeguard the whole; as, for instance, the hand is without deliberation exposed to the blow for the whole body's safety. And since reason copies nature, we find the same inclination among the social virtues; for it behooves the virtuous citizen to expose himself to the danger of death for the public weal of the state\dots{} Consequently, since God is the universal good, and under this good both man and angel and all creatures are comprised, because every creature in regard to its entire being naturally belongs to God, it follows that from natural love angel and man alike love God before themselves and with a greater love. Otherwise, if either of them loved self more than God, it would follow that natural love would be perverse, and that it would not be perfected but destroyed by charity. (St.~Thomas Aquinas, \ST{I, q.~60, a.~5})
\end{quotation}
Self-love exceeds love of another only among coordinate equals; toward the universal Good from whom a creature holds its entire existence, natural love turns outward to God for God's own sake (\ST{I, q.~60, a.~5, ad~1--2}). Natural love still differs from supernatural charity: natural love loves God as universal Author of natural being, whereas charity loves Him as the supernatural object of beatific vision (\ST{I, q.~60, a.~5, ad~4}). Even sinners and demons retain this radical natural inclination toward God as universal good, ``hating'' Him only under particular effects---such as just punishment---that oppose their wicked wills (\ST{I, q.~60, a.~5, ad~5}). Later Scotist and nominalist authors, however, disputed whether pure natural love without charity can genuinely love God above self.

\angeldossier{\ST{I, q.~60, aa.~1--5} (Love or dilection of the angels: natural love, elective love, self-love, mutual natural love, and natural love of God above self).}{Common teaching (aa.~1--5).}{Ecclus~13:19; St.~Augustine, \emph{De Trinitate}~X.1--2 and \emph{De civitate Dei}~XII.9, XIV; Pseudo-Dionysius, \emph{De divinis nominibus}~IV; Aristotle, \emph{Physics}~II and \emph{Nicomachean Ethics}~I.6, III.3, IX.8.}{Several later Scotist and nominalist authors dispute whether an angel or human being can love God above self by pure natural love without supernatural charity (a.~5); none recorded on aa.~1--4.}

\subsection{The Creation of the Angels (\ST{I, q.~61})}

From the angels' nature, intellect, and will, Aquinas turns to their origin in being. First, every angel has an efficient cause of its existence and was created out of nothing by God (\ST{I, q.~61, a.~1}). Psalm~148 summons the heavenly host to praise their Maker: ``Praise ye him, all his angels\dots{} For he spoke, and they were made: he commanded, and they were created'' (Ps~148:2, 5). Because God alone is His own existence, every other being exists only by participation and must derive its entire being from what exists essentially. Against Manichaean dualism and pantheistic emanation, the Church has defined this truth in solemn councils. The Fourth Lateran Council (1215) confesses in \emph{Firmiter credimus} (DS~800) that the one true God is \emph{creator omnium visibilium et invisibilium, spiritualium et corporalium}, ``creator of all things visible and invisible, spiritual and corporeal,'' who \emph{sua omnipotenti virtute simul ab initio temporis utramque de nihilo condidit creaturam, spiritualem et corporalem, angelicam videlicet et mundanam}, ``by His almighty power together from the beginning of time created out of nothing both kinds of creatures, spiritual and corporeal, namely the angelic and the worldly.'' The Council of Florence (\emph{Cantate Domino}, 1442, DS~1333, 1336) and the First Vatican Council (\emph{Dei Filius}, 1870, ch.~1 and cans.~I.1--5) repeat the definition verbatim, anathematizing anyone who claims that spiritual beings emanated from the divine substance or were not produced out of nothing in their whole substance. If Genesis~1 does not name the angels explicitly, St.~Augustine (\emph{City of God}~XI.9) explains that they are signified under ``heaven'' or ``light,'' and Moses refrained from explicit mention lest a people prone to idolatry fall into worshipping incorporeal spirits (\ST{I, q.~61, a.~1, ad~1}).

Second, no angel was produced from eternity (\ST{I, q.~61, a.~2}). Eternal Wisdom declares, ``The Lord possessed me in the beginning of his ways, before he made any thing from the beginning'' (Prov~8:22). Rejecting Peripatetic claims of eternal separate substances, the Catholic Faith defines that God alone---Father, Son, and Holy Ghost---is from eternity, creating all spiritual and material creatures \emph{liberrimo consilio simul ab initio temporis}, ``by a most free counsel, together from the beginning of time'' (Vatican~I, \emph{Dei Filius}, ch.~1). God acts by His free will rather than natural necessity, producing angels after they had not been; although angels stand above the physical time of the heavens, their existence after non-being and the succession of their mental acts mark them as creatures of time (\ST{I, q.~61, a.~2, ad~1--3}).

Third, were the angels created \emph{before} the corporeal world, or \emph{simultaneously} with it (\ST{I, q.~61, a.~3})? Here Aquinas confronts a genuine East-West division among the Fathers with notable charity and restraint. As seen in Section~\ref{sec:patristic-witness}, the great Eastern Fathers---St.~Basil (\emph{Hexaemeron}~I.5), St.~Gregory of Nazianzus (\emph{Orations}~38.9 and~45.5), and St.~John of Damascus (\emph{De fide orthodoxa}~II.3), echoed by St.~Jerome on Titus~1:2---taught that the invisible intellectual cosmos preceded the visible material universe. The Latin tradition, by contrast---led by St.~Augustine (\emph{City of God}~XI.9), Peter Lombard (\emph{Sententiae}~II, d.~2, p.~2, cc.~4--5), and the conciliar wording \emph{simul ab initio temporis}---taught simultaneous creation.

Aquinas judges simultaneous creation to be the more probable opinion, because angels and bodies are not two separate universes, but parts of one universe whose mutual order constitutes the created good; since ``the works of God are perfect'' (Deut~32:4) and no part is perfect separated from its whole, it is fitting that spiritual and corporeal creation began together (\ST{I, q.~61, a.~3}). Yet Aquinas explicitly refuses to condemn the Eastern patristic opinion as erroneous, invoking Jerome's praise of Gregory Nazianzen---``whose authority in Christian doctrine is of such weight that no one has ever raised objection to his teaching, as is also the case with the doctrine of Athanasius''---and showing how defenders of prior creation read ``in the beginning'' in Genesis~1:1 as ``in the Son'' or as the beginning of corporeal time (\ST{I, q.~61, a.~3, ad~3}).

Fourth, Aquinas affirms by fittingness that the angels were created in the empyrean heaven (\ST{I, q.~61, a.~4}). Quoting Walafrid Strabo's \emph{Glossa ordinaria} on Genesis~1:1---that ``heaven'' signifies ``the empyrean, that is, the fiery or intellectual firmament, which is not so styled from its heat, but from its splendor; and which was filled with angels directly it was made'' (a tradition transmitted by Peter Lombard, \emph{Sententiae}~II, d.~2, p.~2, c.~4)---Aquinas emphasizes that angels do not depend on any body for their being, yet creation in the highest corporeal place manifests their governmental presidency over bodily creation (St.~Isidore on Deut~10:14). St.~Augustine's reference to the ``upper atmosphere'' (\emph{De Genesi ad litteram}~III.10) designates either the highest heaven or specifically lower angels assigned to lower bodies, while Lucifer's boast ``I will ascend into heaven'' (Is~14:13) aims not at a physical sky, but at the heaven of the Blessed Trinity (\ST{I, q.~61, a.~4, ad~1--3}).

\angeldossier{\ST{I, q.~61, aa.~1--4} (Production of the angels in natural being: creation \emph{ex nihilo}, beginning in time, simultaneity with the material world, and creation in the empyrean heaven).}{Defined (aa.~1--2); Disputed (a.~3); Thomist position (a.~4).}{Gen~1:1; Deut~10:14, 32:4; Ps~148:2, 5; Prov~8:22; Is~14:13; Fourth Lateran Council, \emph{Firmiter credimus} (DS~800); Council of Florence, \emph{Cantate Domino} (DS~1333, 1336); First Vatican Council, \emph{Dei Filius}, ch.~1 and cans.~I.1--5; St.~Augustine, \emph{De civitate Dei}~XI.9, \emph{De Trinitate}~XIII, and \emph{De Genesi ad litteram}~III.10, VIII.20--21; St.~Jerome, \emph{In Epistolam ad Titum}~1:2; St.~Isidore of Seville; Walafrid Strabo (\emph{Glossa ordinaria} on Gen~1:1); Peter Lombard, \emph{Sententiae}~II, d.~2; \emph{Liber de causis}; Aristotle, \emph{Metaphysics}~VIII.}{Dualist and emanationist errors condemned by Lateran~IV, Florence, and Vatican~I (a.~1); Peripatetic eternity of separate substances (a.~2); the Eastern patristic consensus holding that angels were created prior to the material world---St.~Basil (\emph{Hexaemeron}~I.5), St.~Gregory of Nazianzus (\emph{Oration}~38.9; \emph{Oration}~45.5), and St.~John of Damascus (\emph{De fide orthodoxa}~II.3)---explicitly respected by Aquinas as non-erroneous (a.~3); patristic hesitation over localizing immaterial spirits prior to the corporeal cosmos (a.~4).}

\subsection{Elevation to Grace, Merit, and Beatitude (\ST{I, q.~62})}

Question~62 traces the passage of the holy angels from nature to eternal vision. First, were the angels created already in beatitude (\ST{I, q.~62, a.~1})? Aquinas distinguishes two ultimate perfections. Natural beatitude---the highest contemplation attainable by natural power, which angels possess instantly without discursive search---belonged to them from creation. Supernatural beatitude, however, in which ``we shall see him as he is'' (1~John~3:2), transcends every created nature (a doctrine later protected by the Council of Vienne's decree \emph{Ad nostrum qui} in 1312 against claims that an intellectual nature is naturally blessed without the light of glory). No angel possessed supernatural vision at creation: as St.~Augustine argues (\emph{City of God}~XI.13), genuine beatitude demands unshakeable certainty of eternal perseverance, yet some angels fell.

Second, because eternal life exceeds created capacity---``the grace of God, life everlasting'' (Rom~6:23)---an angel required supernatural grace to turn meritoriously toward God as the object of beatific vision (\ST{I, q.~62, a.~2}). Just as fire cannot generate living flesh unless elevated as an instrument of a living soul, no created will can move toward an end above its nature without a supernatural principle. While angels turned to God naturally as Author of being (\ST{I, q.~60, a.~5}), turning to Him as Bestower of glory demanded grace, without any infinite regress: God moves the will first by preparatory operative grace (Lam~5:21), then through meritorious habitual grace, and finally into consummate glory (\ST{I, q.~62, a.~2, ad~1--3}).

Third, did God bestow sanctifying grace at the very moment of creation, or did He create the angels first in pure natural integrity (\emph{in puris naturalibus}) prior to grace (\ST{I, q.~62, a.~3})? Here the medieval schools genuinely divided. Peter Lombard taught that angelic nature was created formed in natural habits yet still lacking the higher form of charity until conversion, assisted by cooperating grace (\emph{Sententiae}~II, d.~2, p.~2, c.~5; d.~5, c.~4). St.~Bonaventure adopted Lombard's account in \emph{In II Sent.}~(d.~3) and in \emph{Breviloquium}~(II.7), picturing the newly created angels as poised midway between the supreme good and mutable good, ready to ascend into grace and glory if they turned upward. Aquinas records the disagreement openly and judges creation \emph{in sanctifying grace} to be more probable and more faithful to the Saints, citing St.~Augustine's words in \emph{City of God}~(XII.9)---a passage later woven directly into the 1566 \emph{Catechismus Romanus} (Part~I, Art.~1)---that God created the angels with pure love, ``at the same time building up their nature and bestowing grace on them.'' Just as Augustine held that material creation began endowed with seminal principles (\emph{De Genesi ad litteram}~VIII.3), God implanted sanctifying grace---the ``seed'' of God (1~John~3:9)---from the beginning as the principle of merit, inclining the angelic will freely without compelling it (\ST{I, q.~62, a.~3, ad~1--3}).

Fourth and fifth, an angel merited his beatitude (\ST{I, q.~62, a.~4}) and obtained it immediately after a single act of charity (\ST{I, q.~62, a.~5}). Because the heavenly city measures man and angel alike (Apoc~21:17), supernatural glory came to the angels as a reward of virtue rather than an unmerited possession, ruling out both beatitude without merit and any conflation of imperfect merit with consummate fruition in one identical act (\ST{I, q.~62, a.~4}). Yet why did human beings receive a lifelong pilgrimage of trial while angels passed into glory after one single act? Aquinas answers with a foundational principle of Catholic theology: ``grace perfects nature according to the manner of the nature'' (\ST{I, q.~62, a.~5}). Because angelic nature attains its natural perfection at once without gradual stages, an angel attained glory straightway upon one act of charity---meriting in one angelic instant and entering beatific vision in the immediately succeeding instant, just as a separated human soul with merit enters heaven at once unless hindered (Luke~20:36; \ST{I, q.~62, a.~5, ad~1--2}).

Sixth, Aquinas holds with Peter Lombard (\emph{Sententiae}~II, d.~3, p.~2, c.~4) that the angels received grace and glory in proportion to their natural gifts (\ST{I, q.~62, a.~6}). Two reasons converge: on God's side, a wise builder chisels finer stones for the more ornate part of the house, ordaining nobler angelic species for higher degrees of glory; on the angel's side, lacking lower bodily passions that retard or divide the human spirit, each higher angel turned to God with its entire greater natural energy. Aquinas immediately adds a sharp caution (\ST{I, q.~62, a.~6, ad~3}): that proportion does \emph{not} hold among human beings, who differ only numerically within one species and suffer internal resistance from lower appetites. Later Franciscan and Jesuit masters, including John Duns Scotus and Francisco Su\'arez, rejected any necessary proportion between natural angelic rank and degree of glory.

Seventh, beatitude does not erase natural knowledge or natural love (\ST{I, q.~62, a.~7}). Glory perfects nature rather than destroying it; St.~Paul's word that ``when that which is perfect is come, that which is in part shall be done away'' (1~Cor~13:10) removes opposing privations such as faith, not underlying natural perfections, so a blessed angel simultaneously knows God directly in His essence and naturally in the mirror of its own essence (\ST{I, q.~62, a.~7, ad~1--3}).

Eighth, a beatified angel cannot sin (\ST{I, q.~62, a.~8}). Against Origen's hypothesis that beatified wills could grow satiated and fall again---anathematized at Constantinople in 543/553---St.~Augustine affirms that ``there is in the holy angels that nature which cannot sin'' (\emph{De Genesi ad litteram}~XI). Seeing God's essence, which is the very essence of goodness, the beatified angel stands toward God exactly as every rational will stands toward universal goodness: unable to will or act except as ordered to Him. In his celebrated third reply, Aquinas refutes the illusion that sinlessness narrows freedom:
\begin{quotation}
Hence it belongs to the perfection of its liberty for the free-will to be able to choose between opposite things, keeping the order of the end in view; but it comes of the defect of liberty for it to choose anything by turning away from the order of the end; and this is to sin. Hence there is greater liberty of will in the angels, who cannot sin, than there is in ourselves, who can sin. (St.~Thomas Aquinas, \ST{I, q.~62, a.~8, ad~3})
\end{quotation}

Finally, beatified angels do not advance in \emph{essential} beatitude (\ST{I, q.~62, a.~9}). No longer wayfarers (\emph{viatores}) but possessors (\emph{comprehensores}), every blessed spirit is directed by divine predestination to a determinate finite mode of beholding God's infinite essence---one more perfectly than another according to diversity of merits, as the Council of Florence teaches of the blessed (\emph{Laetentur caeli}, DS~1305). Once that final term is secured, essential vision cannot grow, and angelic ministry flows from beatitude rather than meriting it (\ST{I, q.~62, a.~9, ad~1--2}). Yet the angels' \emph{accidental} joy genuinely increases until the Last Judgment as sinners repent through their ministrations---``there shall be joy before the angels of God upon one sinner doing penance'' (Luke~15:10)---whereas Christ alone was simultaneously wayfarer and comprehensor (\ST{I, q.~62, a.~9, ad~3}).

\angeldossier{\ST{I, q.~62, aa.~1--9} (Perfection of the angels in grace and glory: natural versus supernatural beatitude, need for grace, creation in grace, merit, immediacy, proportion to natural gifts, preservation of nature, impeccability, and fixed essential glory).}{Church teaching (aa.~2, 8); Disputed (a.~3); Common teaching (aa.~1, 4, 7, 9); Thomist position (aa.~5--6).}{Lam~5:21; Zach~1:3; Luke~15:10, 20:36; Rom~6:23, 11:6; 1~Cor~13:10; Heb~1:14; 1~John~3:2, 9; Apoc~21:17; Constantinople (543/553) against Origenist mutability; Council of Vienne, \emph{Ad nostrum qui} (1312); Council of Florence, \emph{Laetentur caeli} (DS~1305); \emph{Catechismus Romanus}, Part~I, Art.~1; St.~Augustine, \emph{De civitate Dei}~XI.13, XII.9, \emph{De Genesi ad litteram}~I.15, II.8, IV.34, V, VIII.3, XI, and \emph{De doctrina christiana}~I; Gennadius (\emph{De ecclesiasticis dogmatibus}~XXIX); Peter Lombard, \emph{Sententiae}~II, dd.~3--5; Aristotle, \emph{Nicomachean Ethics}~I.8--9, X and \emph{Metaphysics}~IV, IX.}{Peter Lombard (\emph{Sententiae}~II, d.~2, p.~2, c.~5; d.~3) and St.~Bonaventure (\emph{In II Sent.}~d.~3; \emph{Breviloquium}~II.7) hold creation in natural integrity prior to the infusion of sanctifying grace (a.~3); some medieval authors posited a longer temporal interval (\emph{mora}) of trial before beatification (a.~5); Bl.~John Duns Scotus and Francisco Su\'arez deny a necessary proportion between natural angelic rank and degree of glory (a.~6); Origenist mutability of beatified wills, condemned at Constantinople in 543/553 (a.~8); none recorded on aa.~1--2, 4, 7, 9.}
